\documentclass[10pt,a4paper]{amsart}
\usepackage{amsmath,amssymb,amsthm,graphicx}
% Retained minimal preamble for the Stacks Project chapter sites-cohomology.tex (Cohomology on Sites).
% Verbatim excerpt of the chapter preamble: Xy-pic setup (native preamble lines 14-18) and the native
% theorem declarations with the mathematical macros (native preamble lines 154-217). Copied mechanically;
% nothing added, removed or redefined. The wrapper class and the amsmath/amssymb/amsthm/graphicx base
% packages are supplied by the assembler, so the class/externaldocument part of the native preamble is omitted.
% For commutative diagrams we use Xy-pic
\usepackage[all]{xy}

% We use 2cell for 2-commutative diagrams.
\xyoption{2cell}
% Theorem environments.
%
\theoremstyle{plain}
\newtheorem{theorem}[subsection]{Theorem}
\newtheorem{proposition}[subsection]{Proposition}
\newtheorem{lemma}[subsection]{Lemma}

\theoremstyle{definition}
\newtheorem{definition}[subsection]{Definition}
\newtheorem{example}[subsection]{Example}
\newtheorem{exercise}[subsection]{Exercise}
\newtheorem{situation}[subsection]{Situation}

\theoremstyle{remark}
\newtheorem{remark}[subsection]{Remark}
\newtheorem{remarks}[subsection]{Remarks}

\numberwithin{equation}{subsection}

% Macros
%
\def\lim{\mathop{\mathrm{lim}}\nolimits}
\def\colim{\mathop{\mathrm{colim}}\nolimits}
\def\Spec{\mathop{\mathrm{Spec}}}
\def\Hom{\mathop{\mathrm{Hom}}\nolimits}
\def\Ext{\mathop{\mathrm{Ext}}\nolimits}
\def\SheafHom{\mathop{\mathcal{H}\!\mathit{om}}\nolimits}
\def\SheafExt{\mathop{\mathcal{E}\!\mathit{xt}}\nolimits}
\def\Sch{\mathit{Sch}}
\def\Mor{\mathop{\mathrm{Mor}}\nolimits}
\def\Ob{\mathop{\mathrm{Ob}}\nolimits}
\def\Sh{\mathop{\mathit{Sh}}\nolimits}
\def\NL{\mathop{N\!L}\nolimits}
\def\CH{\mathop{\mathrm{CH}}\nolimits}
\def\proetale{{pro\text{-}\acute{e}tale}}
\def\etale{{\acute{e}tale}}
\def\QCoh{\mathit{QCoh}}
\def\Ker{\mathop{\mathrm{Ker}}}
\def\Im{\mathop{\mathrm{Im}}}
\def\Coker{\mathop{\mathrm{Coker}}}
\def\Coim{\mathop{\mathrm{Coim}}}

% Boxtimes
%
\DeclareMathSymbol{\boxtimes}{\mathbin}{AMSa}{"02}

%
% Macros for moduli stacks/spaces
%
\def\QCohstack{\mathcal{QC}\!\mathit{oh}}
\def\Cohstack{\mathcal{C}\!\mathit{oh}}
\def\Spacesstack{\mathcal{S}\!\mathit{paces}}
\def\Quotfunctor{\mathrm{Quot}}
\def\Hilbfunctor{\mathrm{Hilb}}
\def\Curvesstack{\mathcal{C}\!\mathit{urves}}
\def\Polarizedstack{\mathcal{P}\!\mathit{olarized}}
\def\Complexesstack{\mathcal{C}\!\mathit{omplexes}}
% \Pic is the operator that assigns to X its picard group, usage \Pic(X)
% \Picardstack_{X/B} denotes the Picard stack of X over B
% \Picardfunctor_{X/B} denotes the Picard functor of X over B
\def\Pic{\mathop{\mathrm{Pic}}\nolimits}
\def\Picardstack{\mathcal{P}\!\mathit{ic}}
\def\Picardfunctor{\mathrm{Pic}}
\def\Deformationcategory{\mathcal{D}\!\mathit{ef}}

% --- appended 2026-10-12 (Chief): the project own definitions of the reference/slogan/history
% display environments, copied mechanically from the shared Stacks preamble lines 8-12
% (\usepackage{verbatim} plus the three \newenvironment lines, which the chapter preamble keeps
% verbatim from the project preamble). Nothing added, removed or redefined: these environments
% carry no mathematical content and the project itself typesets them as comments. Without these
% three lines a counted statement whose own text uses them fails with Environment slogan undefined.
\usepackage{verbatim}
\newenvironment{reference}{\comment}{\endcomment}
\newenvironment{slogan}{\comment}{\endcomment}
\newenvironment{history}{\comment}{\endcomment}
\usepackage{cleveref}
\crefname{originalreference}{original reference}{original references}
\Crefname{originalreference}{Original reference}{Original references}
\makeatletter
\newlabel{example-category-to-point}{{08PF}{source}{}{sourcepointer-1}{}}
\newlabel{example-category-to-point@cref}{{[originalreference][0][]08PF}{[1][0][]source}{}{}{}}
\makeatother
\crefname{originalreference}{original reference}{original references}
\Crefname{originalreference}{Original reference}{Original references}
\makeatletter
\newlabel{lemma-compute-by-cosimplicial-resolution}{{08Q9}{source}{}{sourcepointer-2}{}}
\newlabel{lemma-compute-by-cosimplicial-resolution@cref}{{[originalreference][0][]08Q9}{[1][0][]source}{}{}{}}
\makeatother
\crefname{originalreference}{original reference}{original references}
\Crefname{originalreference}{Original reference}{Original references}
\makeatletter
\newlabel{lemma-get-it-now}{{08QA}{source}{}{sourcepointer-3}{}}
\newlabel{lemma-get-it-now@cref}{{[originalreference][0][]08QA}{[1][0][]source}{}{}{}}
\makeatother
\crefname{originalreference}{original reference}{original references}
\Crefname{originalreference}{Original reference}{Original references}
\makeatletter
\newlabel{lemma-product-categories}{{08QB}{source}{}{sourcepointer-4}{}}
\newlabel{lemma-product-categories@cref}{{[originalreference][0][]08QB}{[1][0][]source}{}{}{}}
\makeatother
\crefname{originalreference}{original reference}{original references}
\Crefname{originalreference}{Original reference}{Original references}
\makeatletter
\newlabel{lemma-pullback-tensor-product}{{079U}{source}{}{sourcepointer-5}{}}
\newlabel{lemma-pullback-tensor-product@cref}{{[originalreference][0][]079U}{[1][0][]source}{}{}{}}
\makeatother
\crefname{originallemma}{original lemma}{original lemmas}
\Crefname{originallemma}{Original lemma}{Original lemmas}
\makeatletter
\newlabel{simplicial-lemma-products-homotopy}{{08Q4}{source}{}{sourcepointer-6}{}}
\newlabel{simplicial-lemma-products-homotopy@cref}{{[originallemma][0][]08Q4}{[1][0][]source}{}{}{}}
\makeatother
\title{The lower shriek of a derived tensor product is the tensor product of the lower shrieks}
\author{The Stacks Project Authors}
\date{stacks-project:sites-cohomology.tex:local-snapshot-20261011; GNU FDL 1.2 or later}
\begin{document}
\maketitle
\noindent\textit{Editorial note.} All delivered mathematical text is the author's own, copied from the pinned Stacks Project chapter file sites-cohomology.tex; no statement or proof step was rewritten, reordered or removed.\par\medskip
\noindent\textit{Editorial note.} The counted proof is the author's complete argument: the proof is written out in full: the diagram of categories, the homotopy computation of the morphism sets, the application of the auxiliary lemma, and the final specialisation through the product-categories lemma.\par\medskip
\noindent\textit{Editorial note.} The label lemma-eilenberg-zilber names the counted statement, and the proof environment immediately following it in the source is delivered verbatim as its proof.\par\medskip
\noindent\textit{Editorial note.} The counted statement is carried by the author's own lemma environment in the chapter source.\par\medskip
\noindent\textit{Editorial note.} The statement and proof are copied from the pinned chapter file; the six cited inputs are kept as external pointers with their canonical Stacks tags (the example-labelled pointer is recorded with kind section as the assembler requires).\par\medskip
\noindent\textit{Allowed standard prerequisites (editorial).}\begin{itemize}
\item The category-to-point situation and the functor L pi\_!. Same chapter, the example describing the situation of a category over a point with its sheaves of rings and the functor L pi\_!. Supplies the ambient setting and the functor whose compatibility with the tensor product is asserted.
\item Computation of L pi\_! by a cosimplicial resolution. Same chapter, the lemmas computing L pi\_! through a cosimplicial object and the auxiliary computation for products. Supplies the identifications used for the diagonal and the product situation.
\item Commutation of the pullback with the derived tensor product. Same chapter, the lemma stating that the pullback of a derived tensor product is the tensor product of the pullbacks. Supplies the identity g\textasciicircum{}\{-1\}K = K\_1 tensor\textasciicircum{}L K\_2.
\end{itemize}\medskip
\ifdefined\hypertarget\hypertarget{sourcepointer-1}{}\fi
\noindent\textit{Source reference.} Original reference 08PF: sources/stacks-sites-cohomology/pinned-original-source.tex, source lines 10016--10016. Prerequisite pointer invoked at proof line(s) 10352; it names a standard fact of the same chapter that the counted argument uses, and every step of the counted proof built on it is retained verbatim in the delivered body.\par\medskip
\ifdefined\hypertarget\hypertarget{sourcepointer-2}{}\fi
\noindent\textit{Source reference.} Original reference 08Q9: sources/stacks-sites-cohomology/pinned-original-source.tex, source lines 10229--10229. Prerequisite pointer invoked at proof line(s) 10354; it names a standard fact of the same chapter that the counted argument uses, and every step of the counted proof built on it is retained verbatim in the delivered body.\par\medskip
\ifdefined\hypertarget\hypertarget{sourcepointer-3}{}\fi
\noindent\textit{Source reference.} Original reference 08QA: sources/stacks-sites-cohomology/pinned-original-source.tex, source lines 10298--10298. Prerequisite pointer invoked at proof line(s) 10382; it names a standard fact of the same chapter that the counted argument uses, and every step of the counted proof built on it is retained verbatim in the delivered body.\par\medskip
\ifdefined\hypertarget\hypertarget{sourcepointer-4}{}\fi
\noindent\textit{Source reference.} Original reference 08QB: sources/stacks-sites-cohomology/pinned-original-source.tex, source lines 10316--10316. Prerequisite pointer invoked at proof line(s) 10389; it names a standard fact of the same chapter that the counted argument uses, and every step of the counted proof built on it is retained verbatim in the delivered body.\par\medskip
\ifdefined\hypertarget\hypertarget{sourcepointer-5}{}\fi
\noindent\textit{Source reference.} Original reference 079U: sources/stacks-sites-cohomology/pinned-original-source.tex, source lines 3684--3684. Prerequisite pointer invoked at proof line(s) 10388; it names a standard fact of the same chapter that the counted argument uses, and every step of the counted proof built on it is retained verbatim in the delivered body.\par\medskip
\ifdefined\hypertarget\hypertarget{sourcepointer-6}{}\fi
\noindent\textit{Source reference.} Original lemma 08Q4: sources/stacks-chow/tags.txt, source lines 11047--11047. Prerequisite pointer invoked at proof line(s) 10381; it names a standard fact of a sibling chapter that the counted argument uses, and every step of the counted proof built on it is retained verbatim in the delivered body.\par\medskip
\begin{lemma}
\label{lemma-eilenberg-zilber}
Notation and assumptions as in Example \ref{example-category-to-point}.
If there exists a cosimplicial object $U_\bullet$ of $\mathcal{C}$
such that Lemma \ref{lemma-compute-by-cosimplicial-resolution}
applies, then
$$
L\pi_!(K_1 \otimes^\mathbf{L}_{\underline{B}} K_2) =
L\pi_!(K_1) \otimes^\mathbf{L}_B L\pi_!(K_2)
$$
for all $K_i \in D(\underline{B})$.
\end{lemma}

\begin{proof}
Consider the diagram of categories and functors
$$
\xymatrix{
& & \mathcal{C} \\
\mathcal{C} \ar[r]^-u &
\mathcal{C} \times \mathcal{C} \ar[rd]^{u_2} \ar[ru]_{u_1} \\
& & \mathcal{C}
}
$$
where $u$ is the diagonal functor and $u_i$ are the projection functors.
This gives morphisms of ringed topoi $g$, $g_1$, $g_2$.
For any object $(U_1, U_2)$ of $\mathcal{C}$ we have
$$
\Mor_{\mathcal{C} \times \mathcal{C}}(u(U_\bullet), (U_1, U_2)) =
\Mor_\mathcal{C}(U_\bullet, U_1) \times \Mor_\mathcal{C}(U_\bullet, U_2)
$$
which is homotopy equivalent to a point by
Simplicial, Lemma \ref{simplicial-lemma-products-homotopy}.
Thus Lemma \ref{lemma-get-it-now} gives
$L\pi_!(g^{-1}K) = L(\pi \times \pi)_!(K)$ for any $K$ in
$D(\mathcal{C} \times \mathcal{C}, B)$.
Take $K = g_1^{-1}K_1 \otimes_B^\mathbf{L} g_2^{-1}K_2$.
Then $g^{-1}K = K_1 \otimes^\mathbf{L}_{\underline{B}} K_2$
because $g^{-1} = g^* = Lg^*$ commutes with derived tensor product
(Lemma \ref{lemma-pullback-tensor-product}).
To finish we apply Lemma \ref{lemma-product-categories}.
\end{proof}

\end{document}
